\documentclass[conference]{IEEEtran}
\IEEEoverridecommandlockouts

\usepackage{cite}
\usepackage{epigraph}
\usepackage{amsmath,amssymb,amsfonts}
\usepackage{algorithmic}
\usepackage{graphicx}
\usepackage{textcomp}
\usepackage{hyperref}
\usepackage{xcolor}
\usepackage{lineno}
\usepackage{bbm}
\def\BibTeX{{\rm B\kern-.05em{\sc i\kern-.025em b}\kern-.08em
    T\kern-.1667em\lower.7ex\hbox{E}\kern-.125emX}}
\begin{document}

%% TITLE
\title{Achieving Optimality in Robust Plane Fitting
\\ {\Large \textsc{Exploring Convex Formulations for Problems in Robustness}}
}

%% AUTHOR
\author{

\IEEEauthorblockN{1\textsuperscript{st} Joseph Hobbs}
\IEEEauthorblockA{
\href{mailto:josephrosshobbs@gmail.com}{josephrosshobbs@gmail.com}
}

}

\maketitle

\begin{abstract}
The fitting of an optimal plane to raw 3D points (such as those returned from a LiDAR sensor) is a ubiquitous problem in robotic state estimation and sensor calibration.  Without outliers, the problem is easily reduced to principal component analysis (PCA).  However, in the presence of outliers, existing methods either get trapped in local minima, remain only probabilistically complete, or finish in exponential time in the worst case.  Here, we present a method for determining the globally optimal plane to fit a collection of points under the presence of outliers and in polynomial time.  We first formulate the problem as a mixed-integer quadratic program (MIQP), and then show how an existing relaxation from the literature---dubbed \textit{binary cloning}---converts the MIQP into a convex semidefinite program (SDP) that can be solved to global optimality using off-the-shelf solvers.  We then propose a hybrid RANSAC/SDP algorithm for improved performance relative to the ``vanilla" SDP method, at the cost of probabilistic completeness.  We implement both the PCA-based solver and the SDP-based solver in Python and release the code as open-source (\href{https://github.com/josephrhobbs/robust-plane-fitting.git}{josephrhobbs/robust-plane-fitting}).
\end{abstract}

\begin{IEEEkeywords}
plane fitting, LiDAR, SLAM, convex optimization, semidefinite programming.
\end{IEEEkeywords}

\section{Fitting Planes to Points}

The operation of fitting a plane to a collection of 3D points is a ubiquitous ``primitive" in computer vision \cite{gradients}.  Collections of points, often derived from LiDAR sensors or deep learning-based \textit{monocular depth estimation} (MDE) models, are often corrupted by anisotropic noise and are plagued by \textit{outlier} points, whose presence in the data are purely erroneous.  For a plane with normal vector \( n \in \mathbb{R}^3 \) and offset \( b \in \mathbb{R} \), a generative model for points on a plane is given by

\begin{equation}
	x_i = n(n^\top r_i - b) - r_i + \varepsilon_i
\end{equation}

when \( r_i \in \mathbb{R}^3 \) is a random variable generated by an distribution appropriate to the application and \( \varepsilon \in \mathbb{R}^3 \) is additive noise.  We can modify this generative model to account for \textit{outlier} points (which derive from extraneous objects in the scene, etc.).

\begin{equation}
	x_i = \begin{cases}
		(n^\top r_i - b)n - r_i + \varepsilon_i & \text{inlier} \\
		o_i & \text{outlier}
	\end{cases}
\end{equation}

Here, \( o_i \in \mathbb{R}^3 \) is an arbitrary vector.  A challenging problem in the field of 3D computer vision is \textit{robust plane fitting}.  Given this generative model, how can we recover \( n, b \) in the presence of outliers, and when we do not know which points are outliers?

Many methods exist already for plane fitting in the presence of outliers, including RANSAC \cite{ransac}, robust estimators \cite{gradients}, and Branch-and-Bound \cite{bnb}.  However, each of these faces a unique challenge---RANSAC is only probabilistically complete, robust estimators struggle with local minima, and Branch-and-Bound completes in exponential time (although it is often highly successful empirically).  Here, we explore a \textit{convex} formulation for robust plane fitting, enabling us to find \textit{globally optimal solutions} to the robust plane fitting problem in \textit{polynomial time}.  We also explore a probabilistically complete variant of our algorithm to balance practical performance with strong probabilistic guarantees.

\section{Methods for Robust Plane Fitting}

As previously discussed, many methods currently exist for robust plane fitting.  We group these roughly into three classes: robust estimators \cite{gradients}, RANSAC methods \cite{ransac}, and Branch-and-Bound (BnB) \cite{bnb} methods.

\subsection{Robust Estimation}

Formulations for \textit{robust estimation} involve wrapping residual functions \( r(x) \) with a \textit{robust loss} \( \rho(\cdot) \), where \( \rho: \mathbb{R} \rightarrow \mathbb{R} \) typically\footnote{One notable exception to this is the \textit{Huber loss}, in which the derivative of \( \rho \) converges to a constant nonzero value, rather than zero.} possesses the property \( d \rho/dx \rightarrow 0 \) as \( \left\vert x \right\vert \rightarrow \infty \).  Often, no closed-form estimator exists under a robust loss, but the estimator can be computed easily using ``local" gradient-based solution methods, such as Gauss-Newton \cite{gradients}.  The primary disadvantage of this method is that it can become trapped in local minima, failing to reach the globally optimal (and, most likely, correct) estimator.

\subsection{Random Sample Consensus}

A frequently successful paradigm for robust plane fitting is that of \textit{random sample consensus} (RANSAC).  RANSAC describes a class of \textit{probabilistically complete algorithms}---with each iteration, the probability of finding the global optimum increases, and as the number of iterations grows without limit, this probability converges towards unity.  The fundamental concept behind RANSAC involves selecting a small (and frequently minimal) sample, estimating a global estimator from the small subset, and then determining how many points in the full set are ``in consensus" (inliers according to the algorithm).  The estimator with the most corresponding inliners out of \( N \) attempts is then selected as the final estimator \cite{ransac}.

Though powerful and widely accepted in both industry and research, RANSAC suffers from being only \textit{probabilistically complete}, and has no global optimality \textit{guarantee} for any finite \( N \).  Furthermore, many hundreds or thousands of iterations may be required, which is impractical in real-time or compute-limited applications.

\subsection{Branch-and-Bound Algorithms}

The final paradigm for robust plane fitting we discuss here is that of \textit{branch-and-bound algorithms} (BnB).  BnB guarantees global optimality, unlike both robust estimation and RANSAC.  It works by relaxing mixed-integer constraints to compute lower bounds on the original problem's objective and constructing a ``decision tree" of mixed-integer results \cite{bnb}.  Though in practice BnB is highly successful at a variety of problems, its time complexity is exponential in the worst case.  This is unacceptable for a system that must estimate planes accurately and in real time.

\section{Least-Squares Plane Fitting}

We begin by developing formulations for plane fitting in the absence of outliers, and then we will use these results to guide the development of robust formulations.  The problem of fitting a plane to a collection of observed points \( x \in \mathbb{R}^{3 \times N} \) is naturally expressed as a least-squares optimization problem over a \textit{normal vector} \( n \in \mathbb{R}^3 \) and an \textit{offset} \( b \in \mathbb{R} \).  Least-squares estimation is typically representative of isotropic noise---however, we can easily repeat this derivation using the Mahalanobis distance \( \left\Vert \cdot \right\Vert_\Sigma \) for anisotropic noise, which we leave as future work.  We write the plane normal and plane offset as

\begin{align}
	\min_{n, b} \; & \frac12 \sum_{i = 1}^N \left\Vert n^\top x_i - b \right\Vert_2^2 , \\
	\text{subject to } & n^\top n = 1.
	\label{eq:planefit}
\end{align}

Notice that we impose \( n^\top n = 1 \) to require that the vector \( n \) is a \textit{unit normal vector}, not a vector with arbitrary magnitude.  Here, we explore the standard existing method for (nonrobust) plane fitting: an eigenvalue formulation by reduction to principal component analysis (PCA).  It is not difficult to see that (\ref{eq:planefit}) admits a (nearly) closed-form solution for \( n \).  We first expand the squared \( \ell_2 \) norm to obtain

\begin{align}
	\min_{n, b} \; & \frac12 \sum_{i = 1}^N \left[ n^\top x_i x_i^\top n - 2 b n^\top x_i + b^2 \right] \\
	\text{subject to } & n^\top n = 1.
\end{align}

We can differentiate the objective with respect to \( b \), and as there are no constraints on \( b \), its optimal value is easy to determine.

\begin{equation}
	b^\star = \frac1N \sum_{i=1}^N n^\top x_i
	\label{eq:b}
\end{equation}

Substituting this back into (\ref{eq:planefit}), we get a \textit{quadratically constrained quadratic program} (QCQP) in \( n \) only.

\begin{align}
	\min_{n} \; & \frac12 \sum_{i = 1}^N \left\Vert n^\top x_i - \frac1N n^\top \sum_{j=1}^N x_j \right\Vert_2^2 \\
	\text{subject to } & n^\top n = 1
\end{align}

To refer to the ``centered" point set with its centroid at the origin, we introduce the notation \( a_i := x_i - 1/N \cdot \sum_j x_j \).  With this substitution, the problem reduces even further.

\begin{align}
	\min_{n} \; & \frac12 \sum_{i = 1}^N n^\top a_i a_i^\top n \\
	\text{subject to } & n^\top n = 1
\end{align}

We write as \( H \in \mathbb{R}^{3 \times 3} \) the matrix

\begin{equation*}
	\frac12 \sum_{i = 1}^N a_i a_i^\top
\end{equation*}

to obtain the final QCQP.

\begin{align}
	\min_{n} \; & n^\top H n \\
	\text{subject to } & n^\top n = 1
\end{align}

Despite the problem's nonconvexity---the constraint describes the nonconvex set \( S^2 \)---this is easy to solve to global optimality.  \( n \) is simply the vector associated to \( H \) with \textit{minimum associated eigenvalue}.  As such, we refer to this as an ``eigenvalue formulation", because it reframes the problem in terms of computing the eigenvalues of a matrix in \( \mathbb{R}^{3 \times 3} \).  After computing \( n^\star \), we can of course recover \( b^\star \) using the relation (\ref{eq:b}).

\paragraph*{Note} The astute reader will notice that \( H \) is, up to a constant scaling factor, the sample covariance matrix of the set \( \{ x_i \}_{i=1}^N \).  Computing the eigenvalues and eigenvectors of a sample covariance matrix is a common data science technique called \textit{principal component analysis} (PCA).  Therefore, computing the globally optimal solution (in the least-squares sense) for the normal is equivalent to taking PC3 (``principal component 3") after performing PCA.

\section{Convex Robust Plane Fitting}

\subsection{Formulating the Robust Problem}

As previously discussed, many methods exist for extending this formulation to \textit{robust} plane fitting, when some of the input points are \textit{outlier} points, with no meaning or explanatory power.  Such points must be identified and rejected to ensure the resulting plane explains the remaining \textit{inlier} points effectively.  We will choose to use a \textit{truncated least-squares} approach here, modifying the previous formulation (\ref{eq:planefit}) to use a robust estimator.

\begin{align}
	\min_{n, b} \; & \frac12 \sum_{i = 1}^N \min\left( \left\Vert n^\top x_i - b \right\Vert_2^2, \bar{c}^2 \right) \\
	\text{subject to } & n^\top n = 1
	\label{eq:robust}
\end{align}

Here, the inner \( \min(\cdot, \cdot) \) returns the lesser of its two inputs, and \( \bar{c} \) is an error bound selected by the user.  We direct the user to \cite{quasar} for guidance on selecting \( \bar{c} \).  We see from \cite{quasar} that this can be written as a \textit{mixed-integer quadratic program} (MIQP) by introducing inlier/outlier indicator variables \( \theta_i \) corresponding to each point \( x_i \).  We take \( \theta_i = + 1 \) to indicate that \( x_i \) is an inlier, and we take \( \theta_i = -1 \) to indicate that \( x_i \) is an outlier.

\begin{align}
	\min_{n, b, \theta_i} \; & \frac12 \sum_{i = 1}^N \left[ \frac{1+\theta_i}{2} \left\Vert n^\top x_i - b \right\Vert_2^2 + \frac{1-\theta_i}{2} \bar{c}^2 \right] \\
	\text{subject to } & n^\top n = 1 \\
	& \theta_i \in \{ -1, +1 \}
\end{align}

\subsection{Developing the Quadratic Relaxation}

The mixed-integer constraint \( \theta_i \in \{ -1, +1 \} \) is the only remaining barrier to convexity.  An obvious (though perhaps naive) relaxation is

\begin{equation*}
	\theta_i \in [-1, 1].
\end{equation*}

By removing the mixed-integer constraint, the problem is made convex.  However, in practice we will find that our relaxation is not tight.  Rounding \( \theta_i \) to acceptable values in \( \{ -1, +1 \} \) is a possible fix, but this will likely fail to recover globally optimal solutions.

\subsection{Developing the Semidefinite Relaxation}

An alternative method to the quadratic relaxation is a \textit{semidefinite relaxation} which preserves many of the properties of the original problem.  We apply the \textit{binary cloning} technique introduced in \cite{quasar} (originally used for quaternions) to eliminate \( \theta_i \) from the optimization.  Introducing

\begin{equation}
	n_i := \theta_i n
\end{equation}

we can rewrite the problem in terms of \( n, n_i \) only by observing \( n_i^\top n = \theta_i \).  Notice that we also drop the constant factor for brevity, and impose the constraint \( n n^\top = n_i n_i^\top \) as a substitute for \( \theta_i^2 = 1 \).

\begin{align}
	\min_{n, n_i} \; & \sum_{i = 1}^N \left[ (n^\top a_i)^2 + (n^\top a_i)(n_i^\top a_i) + (1 - n_i^\top n) \bar{c}^2 \right] \\
	\text{subject to } & n^\top n = 1 \\
	& n_i^\top n_i = 1 \\
	& n n^\top = n_i n_i^\top
\end{align}

Expanding these squares, we begin to see more clearly the quadratic formulation.  This is the final step required to implement a semidefinite relaxation.

\begin{align}
	\min_{n, n_i} \; & \sum_{i = 1}^N \left[ a_i^\top n n^\top a_i + a_i^\top n n_i^\top a_i + (1 - n_i^\top n) \bar{c}^2 \right] \\
	\text{subject to } & n^\top n = 1 \\
	& n_i^\top n_i = 1 \\
	& n n^\top = n_i n_i^\top
\end{align}

We introduce the decision variable \( Z \in \mathbb{R}^{(3N + 3) \times (3N + 3)} \) to represent the outer product

\begin{align}
	z &:= \begin{bmatrix}
		n \\
		\vdots \\
		n_i \\
		\vdots
	\end{bmatrix} , \\
	Z &:= z z^\top .
\end{align}

Because the entire objective is linear in the outer products \( n n^\top \) and \( n n_i^\top \), there exists a matrix \( C \) such that the objective reduces to \( \mathrm{tr}(ZC) \).  Furthermore, we rewrite each constraint as an affine constraint on \( Z \).

\begin{align}
	n^\top n = 1 &\Rightarrow \mathrm{tr}([Z]_{n n}) = 1 \\
	n n^\top = n_i n_i^\top &\Rightarrow [Z]_{n n} = [Z]_{n n_i} \\
	n_i^\top n_i = 1 &\Rightarrow \text{redundant} \\
\end{align}

We now \textit{almost} have a convex formulation for the robust plane fitting problem.

\begin{align}
	\min_Z \; & \mathrm{tr}(ZC) \\
	\text{subject to } & \mathrm{tr}([Z]_{n n}) = 1 \\
	& [Z]_{n n} = [Z]_{n n_i} \\
	& \mathrm{rank}(Z) = 1
\end{align}

The final constraint, restricting the rank of \( Z \) to one, is our only limitation to convexity.  Because every rank-one matrix is positive semidefinite, we can relax this constraint to the cone constraint \( Z \succeq 0 \) to make the problem convex.  We also add a few more constraints enforcing symmetry of selected blocks within \( Z \)---this is a slightly tighter relaxation of the rank-one constraint than \( Z \succeq 0 \) alone, and in practice it has been found to improve the tightness of convex relaxations similar to this one \cite{quasar}.

\begin{align}
	\min_Z \; & \mathrm{tr}(ZC) \\
	\text{subject to } & \mathrm{tr}([Z]_{n n}) = 1 \\
	& [Z]_{n n} = [Z]_{n n_i} \\
	& [Z]_{n n_i} = [Z]_{n n_i}^\top \\
	& [Z]_{n_i n_j} = [Z]_{n_i n_j}^\top \; (i \ne j) \\
	& Z \succeq 0
	\label{eq:convex}
\end{align}

Our problem is now convex in \( Z \), and we can solve it easily using ``off-the-shelf" \textit{semidefinite programming} (SDP) software such as MOSEK \cite{mosek}.  A simple \textit{singular value decomposition} (SVD) can recover \( n, n_i \).  Under the convention that singular values are ordered descending, and \( \Sigma := \mathrm{diag}(\sigma_1, \sigma_2, \cdots) \), we have

\begin{align}
	Z^\star &= U \Sigma V^\top , \\
	z^\star &= \begin{bmatrix}
		n^\star \\
		\vdots \\
		n_i^\star \\
		\vdots
	\end{bmatrix} \gets \sqrt{\sigma_1} \cdot U_1 .
\end{align}

We can also recover \( \theta_i^\star = {n^\star}^\top n_i^\star \) if desired.

\section{Improving Efficiency}

Theoretically, due to the convexity of the formulation in (\ref{eq:convex}), we are now able to solve robust plane fitting in polynomial time---a semidefinite solver will solve this in an asymptotic complexity of \( \mathcal{O}(N^6) \).  As the input point set scales to hundreds or even thousands of points, the semidefinite method can quickly become intractable.  Thus, we propose a modified algorithm that trades global optimality for \textit{probabilistic} optimality to a desired confidence and significantly improves asymptotic complexity.

\subsection{Random Sample Consensus, Again}

The RANSAC paradigm suggests an effective method for handling large point sets.  Suppose that, from \( N \) total points, we select \( k \) points according to a uniform distribution, with \( k \) such that we are confident we can fit a plane to global optimality.  We then fit a plane using (\ref{eq:convex}), repeat this process for a total of \( R \) iterations, and return the plane that maximizes the number of predicted inliers \cite{ransac}.

\subsection{Computing Probabilistic Bounds for RANSAC}

Assume we solve the semidefinite program (\ref{eq:convex}) given only the downsampled point set of cardinality \( k < N \).  We also assume that, empirically, \( m < k \) points are required to construct a globally optimal normal estimate.  Our SDP has probability

\begin{equation}
	f := \sum_{i = m}^k {k \choose m} p^m (1 - p)^{(k - m)}
\end{equation}

of correctly rejecting outliers and recovering the correct normal \( n^\star \).  In \( R \) iterations, the probability of a single iteration \textit{not} recovering \( n^\star \) is given by

\begin{equation}
	\mathbb{P}(\text{RANSAC fails}) = (1 - f)^R .
\end{equation}

We fix a \textit{confidence threshold} \( \alpha \) representing the desired probability that RANSAC succeeds (e.g. \( \alpha = 1 - 10^{-4} \)).  We then solve for \( R \).

\begin{align}
	1 - \alpha &= (1 - f)^R \\
	\therefore R &= \left\lceil \frac{\log(1 - \alpha)}{\log(1 - f)} \right\rceil
\end{align}

We note that this variant of RANSAC has time complexity

\begin{equation*}
	\mathcal{O}(R k^6)
\end{equation*}

which is significantly improved over \( \mathcal{O}(N^6) \) when \( k \ll N \) and \( 1 - f \ll 1 \).  Though it requires multiple solutions of (\ref{eq:convex}), we posit that significant downsampling, especially when \( N \) grows to thousands of points, e.g. from a dense LiDAR sensor, will drastically improve practical runtimes of an SDP solver.

\section{Discussion}

We have discussed three standard methods for robust plane fitting (robust estimation, RANSAC, and BnB) and additionally covered a fourth method for robust plane fitting via a semidefinite relaxation of a typically nonconvex MIQP.  In practice, relaxations such as this one are tight up to very high outlier ratios \cite{quasar}.  As convex optimization problems, they converge to globally optimal solutions in polynomial time, which cannot be said about any of the three existing methods.

Semidefinite programming solvers are frequently unavailable for commercial use without a license.  However, we expect existing first-order methods such as operator splitting to outperform second-order methods, such as those implemented in MOSEK, by avoiding the construction of a large \( N^2 \times N^2 \) KKT system \cite{admm, mosek}.  Furthermore, the development of an independent implementation of operator splitting is sufficiently achievable.  We leave the explicit implementation of operator splitting for future work.

\section*{Acknowledgements}

Paul Schroeder, Staff Perception Engineer at Tangram Vision, published a two-part ``Plane Fitting" series on the blog \textit{Tangram Visions} \cite{gradients}.  His publication directly inspired this work.

\section*{Generative AI Statement}

I hold a deep respect for the inspired applied mathematicians, engineers, and researchers working to build the future of generative models for language, code, and images (``generative AI").  However, at the same time, I hold myself to the highest standard of accountability and accuracy in my work.  I also recognize that research, development, and writing are deeply human processes that deserve human attention.  Therefore, it is my guarantee that neither this document nor any associated content (including figures and source code) were created or copyedited by generative AI.

\begin{thebibliography}{00}

\bibitem{gradients} P. Schroeder. ``Plane Fitting, Part 1: A Different Way To Think about Plane Fitting." \textit{Tangram Visions}, Tangram Vision, \href{https://www.tangramvision.com/blog/a-different-way-to-think-about-plane-fitting}{https://www.tangramvision.com/blog/a-different-way-to-think-about-plane-fitting}

\bibitem{ransac} M. A. Fischler and R. C. Bolles. ``Random Sample Consensus: A Paradigm for Model Fitting with Applications to Image Analysis andd Automated Cartography." \textit{Communications of the Association for Computing Machinery}, \href{https://dl.acm.org/doi/10.1145/358669.358692}{https://dl.acm.org/doi/10.1145/358669.358692}

\bibitem{bnb} M. Pilanci and S. Boyd. ``Branch and Bound Methods." Stanford University, \href{https://web.stanford.edu/class/ee364b/lectures/bb_slides.pdf}{https://web.stanford.edu/class/ee364b/lectures/bb\_slides.pdf}

\bibitem{quasar} H. Yang and L. Carlone. ``A Quaternion-based Certifiably Optimal Solution to the Wahba Problem with Outliers." International Conference on Computer Vision, \href{https://arxiv.org/pdf/1905.12536}{https://arxiv.org/pdf/1905.12536}

\bibitem{cvx} CVX Research. ``Software for Disciplined Convex Programming." CVX Research, \href{https://cvxr.com/}{https://cvxr.com/}

\bibitem{mosek} MOSEK ApS. ``MOSEK." MOSEK ApS, \href{https://www.mosek.com/}{https://www.mosek.com/}

\bibitem{cast} L. Shaikewitz, S. Ubellacker, and L. Carlone. ``A Certifiable Algorithm for Simultaneous Shape Estimation and Object Tracking." Laboratory for Information and Decision Systems, Massachusetts Institute of Technology, \href{https://arxiv.org/pdf/2406.16837v1}{https://arxiv.org/pdf/2406.16837v1}

\bibitem{admm} S. Boyd. ``ADMM." Stanford University, \href{https://stanford.edu/~boyd/admm.html}{https://stanford.edu/~boyd/admm.html}

\end{thebibliography}

\end{document}
